%!TEX program = xelatex
\documentclass{article}

% Layout and paper size
\usepackage[a4paper, total={6in, 8in}]{geometry}

% Language
\usepackage[english]{babel}

% Font setup for XeLaTeX
% \usepackage{microtype}
% \usepackage{fontspec}
% \setmainfont{Latin Modern Roman}

% Math packages
\usepackage{amsmath, amssymb, amsfonts}
\usepackage{mathtools}

% Graphics and diagrams
\usepackage{graphicx}
\usepackage{tikz}
\usetikzlibrary{decorations.pathreplacing,arrows.meta}

\input{preamble}


\title{Acoustics of the Gabriel's Horn}
\author{Dias Suleimenov}
\date{\today}

\begin{document}

\maketitle

\section*{Introduction}
One time, I made joke with friends, that I was given task to clean Gabriel's horn out of rust in order to play it. It brings question, how it will actually sound like? So here is the analysis.

\section*{Theoretical setup}
I model the Gabriel's horn with Webster's horn equation (\ref{eq:webster}). This differential equation describes the pressure inside the horn with cross-sectional area $S(x)$ and assumes that waves are approximately plane. For derivation see appendix.

\begin{equation}
    \frac{1}{S(x)} \frac{\partial}{\partial x} \left(S(x) \frac{\partial p}{\partial x}\right) = \frac{1}{c^2} \frac{\partial^2 p}{\partial t^2}
\end{equation}
Where, $p(x, t)$ is the pressure in the horn and $c$ is speed of the sound.

Substituting cross-sectional area of horn $S(x) = \frac{\pi}{x^2}$, we get following equation.

\begin{equation}
    \frac{x^2}{\pi} \frac{\partial}{\partial x} \left(\frac{\pi}{x^2}  \frac{\partial p}{\partial x}\right) = \frac{1}{c^2} \frac{\partial^2 p}{\partial t^2}
\end{equation}

% NOT SURE, Analyze it well. 
Now we need focus on waves that exhibit the harmonic oscilations, so we make substitution of $p(x, t) = \psi(x) \cdot e^{-i \omega t}$.
\begin{equation}
    \frac{x^2}{\pi} \frac{\partial}{\partial x} \left(\frac{\pi}{x^2}  \frac{\partial \psi}{\partial x} e^{-i \omega t}\right) = - \left(\frac{\omega}{c}\right)^2 \psi(x).
\end{equation}
We introduce new variable $k = w/c$, for simplification of notation. 

\dots

\begin{equation}
    x^2 \left( \frac{\partial}{\partial x} \left( \frac{1}{x^2} \right) \frac{\partial \psi}{\partial x} + \frac{1}{x^2} \frac{\partial^2 \psi}{\partial x^2}\right)
\end{equation}

\dots

\begin{equation}
    \frac{\partial^2 \psi}{\partial x^2} - \frac{2}{x} \frac{\partial \psi}{\partial x} + k^2 \psi = 0.
\end{equation}

We impose initional condition such as that at throat ($x = 1$) is open end, $\frac{\partial \psi}{\partial x} = 0$, and closed end at $x = L$ with $\psi{L, t} = 0$.

The system that we described can be solved by Bessel equations, with some \dots

We arrive at 

\begin{equation}
    k (L - 1) + \arctan(k) = n \pi
\end{equation}

We solve following using fixed point iterations algorithm.


% Appendix: ADD derivation of equation
\end{document}